\chapter{Conclusão}
\chapter{Considerações finais}
\label{capitulo5}

Uma dimensão atravessa os cinco registros analíticos e pede ênfase antes do fechamento do capítulo: os mecanismos documentados aqui agem através dos arranjos que governaram (e deixaram de governar) o vínculo. Isso significa que a análise se dirige aos arranjos, não às pessoas que os ocuparam. Fábio e Cláudio demonstraram, ao longo das entrevistas, reflexividade sobre tensões, conhecimento dos ciclos de negócio da \texttt{IBM} e compromisso com excelência científica e interesse público. Quando Cláudio descreveu o projeto de línguas indígenas como exemplo de \textit{AI for Good} que \enquote{vai virar algum dia produto... jamais. E não se espera que vire} (PINHANEZ, 2025, informação verbal), ele articulou uma valorização da pesquisa com impacto social irredutível à lógica comercial. Quando Fábio insistiu que \enquote{a gente não faz nenhum projeto para a IBM} (COZMAN, 2025, informação verbal), articulou sua experiência concreta de autonomia acadêmica, confirmada pelo fato de o \texttt{C4AI} produzir, ao longo de cinco anos, \textit{datasets} abertos, publicações em conferências internacionais e ferramentas disponibilizadas sob licenças permissivas.

A coexistência entre essa autonomia declarada e as reorientações documentadas da agenda não configura contradição. Partes em qualquer colaboração influenciam mutuamente suas prioridades, e, quando a translação funciona, essa influência é experimentada como colaborativa mais do que como coerção \parencite{Latour2011CienciaEmAcao}. A dificuldade é que a translação bem-sucedida tende ao que Latour chama de \textit{black-boxing}, isto é, as negociações e os diferenciais de poder que deram forma ao arranjo tornam-se invisíveis. A ausência de arranjos explícitos de governança no \texttt{C4AI} significou que essas dinâmicas de translação se deram informalmente, sem procedimentos documentados que tornassem rastreável o processo decisório. Sem mecanismos institucionais que garantissem a rastreabilidade, é impossível verificar se os interesses foram de fato negociados ou se o alinhamento refletiu capacidade assimétrica de definir o benefício compartilhado.

Callon, Lascoumes e Barthe formulam, em sua análise de \enquote{democracia técnica}, que as assimetrias consequentes raramente decorrem da conduta de atores individuais \parencite{Callon2011Acting}. Elas decorrem de arranjos sociotécnicos que distribuem de forma desigual capacidades, recursos e poder de definição das agendas. No caso do \texttt{C4AI}, o que estreitava de antemão as escolhas possíveis inclui o subfinanciamento crônico da pesquisa pública brasileira, a dependência de recursos corporativos para infraestrutura computacional considerada básica, as pressões institucionais por \textit{impacto} e \textit{inovação} definidos em termos que valem também em registros corporativos, e as assimetrias geopolíticas que fazem da mão de obra qualificada brasileira, na formulação de Cláudio, \enquote{barata demais} em comparações internacionais. Neste espaço, as decisões de aceitar ou recusar financiamento corporativo, publicar em \textit{open access} ou em periódicos fechados, colaborar ou não com pesquisadores \texttt{IBM}, ocorrem dentro de margens estreitas. Como Fábio observou ao explicar por que projetos em saúde e agronegócio permaneceram ativos mesmo após a \texttt{IBM} reorientar os interesses: \enquote{a gente sempre mantém uma certa independência porque a gente tem recursos da Fapesp} (COZMAN, 2025, informação verbal). Projetos que deixaram de ser de interesse da \texttt{IBM} continuaram com recursos da \texttt{FAPESP}. Essa \enquote{certa independência} parcial e constantemente renegociada é o espaço real de agência disponível para pesquisadores em contextos periféricos do Sul Global.

Ao documentar essas dinâmicas, este capítulo não propõe que arranjos público-privados sejam intrinsecamente problemáticos nem que toda colaboração com corporações configure captura disfarçada. Os centros derivados, os pesquisadores formados e os \textit{datasets} abertos são contribuições verificáveis para a capacidade brasileira em IA. Os cinco registros indicam, contudo, que acordos apresentados como horizontais e mutuamente benéficos podem obscurecer assimetrias estruturais de poder, de recursos e de capacidade de definir os termos da colaboração quando a governança permanece informal. O \texttt{C4AI}, nesse registro, é uma existência parcial no sentido em que a Lição~2 do Capítulo~\ref{capitulo1} formula a partir de Strathern: conectado simultaneamente à autonomia acadêmica e à racionalidade corporativa sem ser redutível a nenhuma das duas. A \texttt{IBM} foi, ao mesmo tempo, constitutiva e periférica à pesquisa que financiou. Fábio e Cláudio foram, ao mesmo tempo, atores da universidade pública e porta-vozes de um arranjo corporativo. Os grupos de pesquisa foram ao mesmo tempo do \texttt{C4AI} e não eram, como mostra a trajetória dos centros derivados que partilharam a origem e deixaram de compartilhar o vínculo. É também aqui que se resolve a tensão que anunciei na abertura, quando disse que tomar a \texttt{IBM} como objeto experimentava, na prática, o esforço de fazer alianças sem ferir os sentimentos estabelecidos, no sentido de Stengers. Dirigir a análise aos arranjos, e não às pessoas que os ocuparam, é o que me permitiu descrever criticamente a rede da qual eu também era parte sem converter Fábio e Cláudio em réus: a aliança que mantive com eles e com o centro não se desfez ao descrever as assimetrias, porque a crítica recaiu sobre os dispositivos de governança que faltaram, e não sobre quem agiu dentro das margens estreitas que esses dispositivos deixavam. Foi esse o modo que encontrei de honrar a confiança que me foi dada sem abrir mão da descrição que me propus a fazer.

Recolho esses fios. Este capítulo documentou um arranjo público-privado em inteligência artificial desde a constituição em outubro de 2020 até o encerramento em dezembro de 2025, ao enredar dois horizontes temporais. No horizonte recente, a etnografia acompanhou quase quatro anos de funcionamento do \texttt{C4AI} como acordo entre \texttt{USP}, \texttt{FAPESP} e \texttt{IBM}, até a comunicação interna de não-renovação em agosto de 2025 e a dissolução pública entre novembro e dezembro do mesmo ano. No horizonte de mais de um século, a trajetória rastreou 135 anos de estratégias \texttt{IBM} de construção de dependências tecnológicas, da tabulação de Hollerith no Censo dos EUA de 1890 aos ecossistemas contemporâneos de IA cognitiva. Os cinco registros analíticos que sistematizei, a assimetria orçamentária, a multiplicação institucional, a reorientação sistemática da agenda, a seletividade geográfica e o arranjo que não se inscreveu, descrevem padrões observáveis do arranjo que emergem do seguir dos atores.

A incubação bem-sucedida do \texttt{C4AI}, em sua ambivalência, atravessa todo o caso. Pelo que se costuma contar como sucesso, o \texttt{C4AI} funcionou: múltiplos centros catalisados, \textit{datasets} liberados sob licenças permissivas, centenas de estudantes formados, publicações em conferências internacionais, reconhecimento no terceiro Fórum Global da UNESCO sobre Ética em Inteligência Artificial em Bangkok, em junho de 2025 \parencite{C4AI2025}. O arranjo, ainda assim, foi encerrado. A multiplicação que o sustentava criou as próprias condições do encerramento: os centros derivados asseguraram financiamento autônomo, e, a partir desse ponto, o incentivo corporativo para sustentar o investimento original diminuiu. De uma perspectiva, a \enquote{missão} se cumpriu: o acordo construiu capacidade duradoura. De outra, a mesma evidência põe em jogo a sustentabilidade do vínculo e a ausência de planejamento de transições. As duas leituras incidem sobre a mesma cadeia de associações que o capítulo rastreou.

A clareza maior que este capítulo entrega está naquilo que os arranjos não registravam. A \texttt{IBM} seguiu as provisões contratuais existentes e estava formalmente no direito de não renovar. O que permanece é que as provisões nunca responderam às perguntas elementares sobre o propósito da colaboração, os critérios para sua continuação e os procedimentos de transição. A parte pública comprometeu-se com desenvolvimento institucional de longo prazo, ao passo que a parte corporativa reteve flexibilidade para ajustar compromissos. Pesquisadores brasileiros em regime de bolsa, trabalhando fora das proteções laborais convencionais, enfrentaram dificuldades particulares durante a transição. Essas dinâmicas produzem efeitos que arranjos de governança formalizados teriam condições de antecipar.

A articulação da estratégia de \enquote{ecossistema} que observei na \textit{bluetalks} indica que os porta-vozes corporativos agiam com clareza sobre a posição que ocupavam. A parte pública, ao longo dos cinco anos, moveu-se dentro de um arranjo cujos dispositivos escritos não registravam os termos que o encerramento viria a tornar urgentes. A descrição desses registros de ausência, o descompasso entre horizontes temporais, a concentração de fontes, a invisibilidade do processo decisório sobre reorientações, a condição dos bolsistas, a dependência infraestrutural e epistêmica, a opacidade da consulta estratégica, foi reunida na seção~\ref{subsec:arranjo_nao_inscreveu}.

A seletividade geográfica tensiona os limites do que a etnografia pode estabelecer. A \texttt{IBM} concluiu suas atividades no Brasil simultaneamente à manutenção e à expansão de iniciativas na Índia (Hybrid Cloud Lab com o IISc em 2021, Quantum Valley Tech Park com a TCS e o governo de Andhra Pradesh em maio de 2025, BharatGen em setembro de 2025) e nos EUA (Genesis Mission em novembro de 2025). Fatores comparativos de escala, tamanho de mercado e prioridades estratégicas provavelmente influenciaram essas decisões, e os pesos específicos permanecem indisponíveis. O que o arranjo não inscreveu aparece aqui como opacidade, condição de funcionamento que atravessa a confiança e o planejamento de todas as partes.

O que descrevi aqui, no plano institucional, ainda não responde por completo à pergunta etnográfica que levantei no Capítulo~\ref{capitulo1}: onde e como pesquisadores fazem IA no cotidiano? Como se materializaram, durante os cinco anos de funcionamento do \texttt{C4AI}, as dependências infraestruturais e as assimetrias geopolíticas que identifiquei aqui? Como os pesquisadores negociaram algoritmos, dados, infraestruturas computacionais limitadas e demandas institucionais em tensão? O próximo capítulo examina essas perguntas pelo sítio etnográfico do desenvolvimento do \texttt{Spira} durante a pandemia de Covid-19, um dos projetos ativos durante o período de vigência do arranjo. Acompanho Marcelo Finger e os múltiplos actantes que se reuniram na construção de IA preditiva em contexto de urgência sanitária. Trata-se de mudar de ponto de entrada, porque a rede não tem topo nem fundo, e a passagem ao próximo capítulo não desce de um nível superior a um inferior: parto dos arranjos institucionais e das trajetórias históricas para as práticas cotidianas e as translações situadas, outro ponto da mesma rede a partir do qual outras associações se tornam seguíveis. O que aparece, desse outro ponto, é que fazer inteligência artificial no Brasil daquele período foi trabalho situado, material e político, atravessado pelas mesmas dependências estruturais que identifiquei aqui, e negociado de formas criativas pelos pesquisadores no dia a dia do laboratório.

Em dezembro de 2025, quando encerrei a redação deste capítulo, a rede estava em reconfiguração. O que permanecia em jogo para os actantes que restavam no centro era a redistribuição dos recursos, das competências e dos vínculos institucionais que o arranjo havia organizado por cinco anos: o financiamento da \texttt{FAPESP}, a infraestrutura computacional do JAIRU recém-inaugurado, os grupos que haviam conquistado autonomia e os que permaneciam dependentes dos recursos conjuntos, e as relações com as corporações que poderiam ou não suceder o papel exercido pela \texttt{IBM}. Nenhum desses elementos estava definido no momento em que esta escrita se encerrou. As conversas informais com Fábio e Cláudio naquele período registravam interpretações ainda abertas, antes que a narrativa definitiva sobre o que havia acontecido se fechasse. A encruzilhada, nos termos de Latour, é o ponto em que os porta-vozes ainda não sabem com quantos atores falam \parencite{Latour2011CienciaEmAcao}. É exatamente aí que esta etnografia encerrou o registro.

Volto, antes de encerrar, às perguntas com que abri este capítulo, para dizer o que o seguir dos atores me deixou afirmar sobre cada uma, sem fechá-las em respostas. Perguntei por quais modos uma corporação passa a extrair valor de suas conexões com universidades: o que pude descrever foi um conjunto de modos rastreáveis, a aprendizagem barata que Cláudio nomeou, a multiplicação de centros que a atividade do \texttt{C4AI} catalisou, a reorientação da agenda conforme as frentes corporativas se deslocavam, sem que eu possa fixar, a partir do que segui, a intenção que teria guiado cada movimento. Perguntei o que significa \enquote{fortalecer um ecossistema de IA}: vi esse fortalecimento e a atuação comercial da própria \texttt{IBM} correrem juntos, e descrevi essa coincidência sem precisar atribuir à empresa o cálculo de fazer um gerar o outro. Perguntei de que maneira se \enquote{descobre} um modelo de negócios baseado em código aberto e como essa descoberta orienta políticas globais: rastreei a fala de Cláudio sobre os anos 1990 e a aquisição da Red Hat até uma trajetória mais longa, em que a abertura técnica e o fechamento infraestrutural se sustentam mutuamente, sem que a fala comprove a trajetória nem a trajetória comprove a fala. E perguntei como tudo isso se relaciona com um centro de IA numa universidade pública brasileira: o \texttt{C4AI} apareceu como o ponto em que essas racionalidades se inscreveram localmente, sob a forma de um arranjo que se compôs com facilidade e se desfez com a mesma facilidade que o instrumento de 2016 já previa. As perguntas que permaneceram abertas não se fecharam por falta de rastreamento; elas tocam no que corre por canais que não acompanhei, e é dessa fronteira que trato a seguir.

Retomo, para encerrar, a pergunta que a \textit{bluetalks} me deu e que tomei como questão deste capítulo: o que circula nessas cadeias, por quais caminhos e sob que termos. A circulação de valor que pude descrever é a que passa pelos porta-vozes que segui e pelos documentos que pude consultar. O que corre por canais privados, os cálculos internos, as negociações bilaterais e as deliberações estratégicas, não se ofereceu à descrição, e é esse o ponto que prefiro não suavizar: a opacidade do núcleo decisório é parte do que circula nessas cadeias e condição de seu funcionamento, e não uma falha do meu rastreamento. Foi o corte que fiz ao seguir os porta-vozes, no sentido de \textcite{strathern2011cortando}, que me deu a circulação institucional e me subtraiu a decisória, e descrever esse corte como tal, em vez de simular um acesso que não tive, é o que mantém a descrição honesta quanto ao que a minha posição alcança. Segui os atores como Latour recomenda \parencite{Latour2011CienciaEmAcao} e documentei aquilo que escolheram tornar visível, ciente de que o que permanece invisível molda os desfechos de forma substantiva. O que este capítulo entrega, então, é a descrição do que age informalmente nos arranjos que produziu, e da visibilidade que o encerramento lançou sobre essa informalidade. O que o caso não resolve, e que a evidência etnográfica por si só não tem condições de resolver, é uma pergunta sobre o papel que resta à universidade quando a fronteira da produção de conhecimento em IA migra para a infraestrutura corporativa. Se os vínculos com quem comanda essa infraestrutura se tornaram condição de funcionamento para a pesquisa brasileira em IA, a pergunta é sob quais arranjos eles vêm sendo buscados, com quais salvaguardas, e em direção a quais finalidades. A essa pergunta não entrego uma conclusão; ofereço um problema que a descrição tornou visível. É esse problema que o próximo capítulo, ao descer para o cotidiano do \texttt{Spira}, ajuda a recolocar desde dentro do laboratório.
Retomo, para encerrar, a pergunta que a \textit{bluetalks} me deu e que tomei como questão deste capítulo: o que circula nessas cadeias, por quais caminhos e sob que termos. A circulação de valor que pude descrever é a que passa pelos porta-vozes que segui e pelos documentos que pude consultar. O que corre por canais privados, os cálculos internos, as negociações bilaterais e as deliberações estratégicas, não se ofereceu à descrição, e é esse o ponto que prefiro não suavizar: a opacidade do núcleo decisório é parte do que circula nessas cadeias e condição de seu funcionamento, e não uma falha do meu rastreamento. Foi o corte que fiz ao seguir os porta-vozes, no sentido de \textcite{strathern2011cortando}, que me deu a circulação institucional e me subtraiu a decisória, e descrever esse corte como tal, em vez de simular um acesso que não tive, é o que mantém a descrição honesta quanto ao que a minha posição alcança. Segui os atores como Latour recomenda \parencite{Latour2011CienciaEmAcao} e documentei aquilo que escolheram tornar visível, ciente de que o que permanece invisível molda os desfechos de forma substantiva. O que este capítulo entrega, então, é a descrição do que age informalmente nos arranjos que produziu, e da visibilidade que o encerramento lançou sobre essa informalidade. O que o caso não resolve, e que a evidência etnográfica por si só não tem condições de resolver, é uma pergunta sobre o papel que resta à universidade quando a fronteira da produção de conhecimento em IA migra para a infraestrutura corporativa. Se os vínculos com quem comanda essa infraestrutura se tornaram condição de funcionamento para a pesquisa brasileira em IA, a pergunta é sob quais arranjos eles vêm sendo buscados, com quais salvaguardas, e em direção a quais finalidades. A essa pergunta não entrego uma conclusão; ofereço um problema que a descrição tornou visível. É esse problema que o próximo capítulo, ao descer para o cotidiano do \texttt{Spira}, ajuda a recolocar desde dentro do laboratório.
